\scanpage{22}
\mspage{2}
Je prouve dans \unsure{papiers bleus}

\uline{Th} \uline{C'est un isom} ($X$ \deleted{\unsure{pré}}schéma loc. de type fini \deleted{séparé} sur $\mathbb C$)

\inserted{Si $X$ affine, cela équivaut au :}
\uline{Corollaire} \deleted{\unsure{immédiat|Théorème}}
Si $X$ affine non sing., alors $H^*(X^{\mathrm{an}},\mathbb C)$ peut se calculer avec le complexe de De Rham ``algébrique''

N.B. Le théorème résulte facilement de ce cas particulier, par suite spectrale de Leray d'un recouvrement.

Indications sur démonstration. Si $X$ est \uline{propre}, cela résulte des suites spectrales de type *), algébrique, et analytique, et de Gaga, impliquant qu'on a un isom sur les termes initiaux. [C'est donc essentiellement connu depuis Gaga de Serre]

Le cas général, donc \unsure{m$^{\mathrm e}$|m$^{\mathrm t}$} le cas affine, demande un moyen plus puissant que Gaga, la résolution des singularités. Donner \deleted{\unsure{par}} une idée du principe de la démonstration.

4) \uline{Conséquences métaphysiques} \inserted{\uline{du th. de comparaison}.} La cohomologie complexe (définie par voie transcendante), interprétée comme ``cohomologie de De Rham''
\margin{21}
